% TOP_PROBLEM: {"problemId": "problem.the-nc-versus-p-problem-for-efficient-parallel-computation", "canonicalTitle": "The NC versus P Problem for Efficient Parallel Computation", "releaseRank": 23, "primaryDomain": "theoretical_computer_science", "primaryDomainLabel": "Theoretical computer science", "displayStatus": "open", "releaseStatus": "open"}
% !TeX program = lualatex
% Definition and sources only; this file is not an LLM solution attempt.
\documentclass[11pt]{article}
\usepackage[margin=25mm]{geometry}
\usepackage{fontspec}
\setmainfont{Latin Modern Roman}
\usepackage{amsmath,amssymb,mathtools}
\usepackage{unicode-math}
\setmathfont{Latin Modern Math}
\usepackage{xurl,hyperref}
\hypersetup{colorlinks=true,urlcolor=blue}
\setcounter{secnumdepth}{0}
\setlength{\parindent}{0pt}
\setlength{\parskip}{5pt}
\setlength{\emergencystretch}{3em}
\begin{document}
\section*{\texorpdfstring{The NC versus P Problem for Efficient Parallel Computation}{The NC versus P Problem for Efficient Parallel Computation}}\label{23}
\subsection{Definitions and mathematical statement}
For bounded-fan-in Boolean circuits, size is the number of gates and depth is the longest input-to-output path. In the standard uniform convention, $\mathsf{NC}=\bigcup_{k\ge1}\mathsf{NC}^k$, where $\mathsf{NC}^k$ has polynomial size and depth $O((\log n)^k)$. Uniformity requires an effective resource-bounded procedure generating the circuit family, rather than an arbitrary circuit for each length. The question is
\[
 \mathsf{NC}\stackrel{?}{=}\mathsf P.
\]
For precision, the conventional logspace-uniform interpretation is used; the catalogue does not choose among finer uniformity conventions.
\par\smallskip\textit{Formulation and concept references:} \href{https://arxiv.org/abs/2508.20735}{[S1]}.

\subsection{Short English statement}
Can every efficient sequential decision algorithm be replaced by a highly parallel one using only polylogarithmically many stages?
\subsection{Sources}
\begin{itemize}
\item[S1]J. Heemstra, J. Martens, and A. Wijs, Evaluating Massively Parallel Algorithms for DFA Minimisation, Equivalence Checking and Inclusion Checking, Introduction (2026).
\url{https://arxiv.org/abs/2508.20735}.
\textit{Formulation source.}
\item[S2]Parallel Computation, Matt Williamson.
\url{https://community.wvu.edu/~krsubramani/courses/sp09/cc/lecnotes/parallel.pdf}.
\textit{Further reference; not a proof endorsement.}
\item[S3]Complexity of the Freezing Majority Rule with L-shaped Neighborhoods.
\url{https://arxiv.org/html/2509.16065v1}.
\textit{Further reference; not a proof endorsement.}
\end{itemize}
\end{document}
